%%%%%%%%%%%%%%%%
%%% gen 21 %%%
%%%%%%%%%%%%%%%%

\Chapter{21}

\Verse{1}
{
	\hJ{וַיהֹוָה פָּקַד אֶת שָׂרָה כַּאֲשֶׁר אָמָר}
	\hP{וַיַּעַשׂ יְהֹוָה לְשָׂרָה כַּאֲשֶׁר דִּבֵּר׃}
}
{
	\eJ{And YHWH took account of Sarah as He had said,}
	\eP{and YHWH did to Sarah as He had spoken.}
}
{
	YHWH gets Sarah pregnant like he promised.
}

\Verse{2}
{
	\hJPP{וַתַּהַר וַתֵּלֶד שָׂרָה לְאַבְרָהָם בֵּן לִזְקֻנָיו לַמּוֹעֵד אֲשֶׁר דִּבֶּר אֹתוֹ אֱלֹהִים׃}
}
{
	\eJPP{
		And Sarah became pregnant and gave birth to a son for
		Abraham in his old age at the appointed time that God had
		spoken.
	}
}
{
	She has the baby.

	Because \eRed{God} got her pregnant \eP{at the appointed time}.\eRed{\footnote{
		Documentary Note: This line is often attributed to J.
		This is strange, given that it contains the word ``God'' spoken by the narrator.
		This is otherwise unknown in the J source, as J herself never calls
		YHWH ``God,'' though characters in J sometimes do (the first example being the
		snake in Genesis 3.) Also the use of the phrase ``appointed time'' is
		characteristic of P, though not exclusively.
	}}
}

\Verse{3}
{
	\hP{וַיִּקְרָא אַבְרָהָם אֶת שֶׁם בְּנוֹ הַנּוֹלַד לוֹ אֲשֶׁר יָלְדָה לּוֹ שָׂרָה יִצְחָק׃}
}
{
	\eP{
		And Abraham called the name of his son who was born to
		him, to whom Sarah had given birth for him: Isaac.
	}
}
{
	So they name him Laughter.
}

\Verse{4}
{
	\hP{וַיָּמל אַבְרָהָם אֶת יִצְחָק בְּנוֹ בֶּן שְׁמֹנַת יָמִים כַּאֲשֶׁר צִוָּה אֹתוֹ אֱלֹהִים׃}
}
{
	\eP{
		And Abraham circumcised Isaac, his son, at eight days old,
		as God had commanded him.
	}
}
{
	He cuts the covenant on day eight.
}

\Verse{5}
{
	\hP{וְאַבְרָהָם בֶּן מְאַת שָׁנָה בְּהִוָּלֶד לוֹ אֵת יִצְחָק בְּנוֹ׃}
}
{
	\eP{
		And Abraham was a hundred years old when Isaac, his son,
		was born to him.
	}
}
{
	Now Abe's a hundred.
}

\Verse{6}
{
	\hE{וַתֹּאמֶר שָׂרָה צְחֹק עָשָׂה לִי אֱלֹהִים כּל הַשֹּׁמֵעַ יִצְחַק לִי׃}
}
{
	\eE{
		And Sarah said, “God has made laughter for me. Everyone
		who hears will laugh for me.”
	}
}
{
	Sarah speaks to no one in particular, reminding the audience her son's a pun.
}

\Verse{7}
{
	\hJ{וַתֹּאמֶר מִי מִלֵּל לְאַבְרָהָם הֵינִיקָה בָנִים שָׂרָה כִּי יָלַדְתִּי בֵן לִזְקֻנָיו׃}
}
{
	\eJ{
		And she said, “Who would have said to Abraham, ‘Sarah has
		nursed children’? Yet I’ve given birth to a son in his old
		age.”
	}
}
{
	Now Sarah speaks

	She's like ``Who would've thought this would happen eh?''
}

\Verse{8}
{
	\hE{וַיִּגְדַּל הַיֶּלֶד וַיִּגָּמַל וַיַּעַשׂ אַבְרָהָם מִשְׁתֶּה גָדוֹל בְּיוֹם הִגָּמֵל אֶת יִצְחָק׃}
}
{
	\eE{
		And the boy grew and was weaned. And Abraham made a big
		feast in the day that Isaac was weaned.
	}
}
{
	He grows up and stops nursing. Abe makes a big dinner.
}

\Verse{9}
{
	And Sarah looks over...

	\hE{וַתֵּרֶא שָׂרָה אֶת בֶּן הָגָר הַמִּצְרִית אֲשֶׁר יָלְדָה לְאַבְרָהָם מְצַחֵק׃}
}
{
	\eE{
		And Sarah saw the son of Hagar, the Egyptian, whom she had
		borne to Abraham, fooling around.
	}
}
{
	Sarah looks over at sisterwife's son.

	The text says she sees him ``fooling around,'' which is a pun on Isaac. (No time to explain, just trust me, the Hebrew word's not the point here.)

	\aC{The Greek text has “fooling around with her son Isaac,” which can connote either playing with Isaac or mocking him. The word in Hebrew is \heb{מצחק}, which is a pun on the name Isaac. It will occur again in the story of Isaac, Rebekah, and Abimelek (Gen 26:8).}

	Anyways the point is that this new E character seems to like puns too.

	To recap:

	\begin{itemize}
	\item J's entire universe is stitched together from a dense fabric of puns and sex and subtext.
	\item E's universe is one in which God exists and his name's YHWH\footnote{Though in E we don't learn this until Exodus 3.} but it's not illegal to have fun, and in which writing a bible doesn't rule out the occasional pun.
	\item P's universe is one in which language, all languages, and each language according to its words, are a tool of holiness. It is a holy thing. And in the holy medium of language, all languages, and each language according to its words, puns are an abomination and simply must not be done. When anyone brings a pun into the world by bringing it into one's mind and conceiving of it, he must bring his pun or puns to the Tabernacle and present it in front of Aaron and his sons. And he shall lay his hand on the head of the pun and slaughter it at the entrance of the tent of meeting. And Aaron's sons the priests shall sprinkle the blood all around against the sides of the altar. Then he shall remove all the letters from the insides of the word of his pun (if it is a single word pun) or the words of his pun (if it is a multiword pun), and then Aaron's sons are to burn it to smoke on the altar; it is a literative offering, a literation pleasing to the Lord. After these things, any male in a priest’s family may hear the words of the pun; it is a holy thing. I am YHWH who makes you holy, who brought you out of the land of Egypt, to be your God. I am YHWH.
	\end{itemize}

%	\footnote{
%		fooling around. The Greek text has “fooling around with her son Isaac,”
%		which can connote either playing with Isaac or mocking him. The word in
%		Hebrew is , which is a pun on the name Isaac. It will occur again in the
%		story of Isaac, Re-bekah, and Abimelek (Gen 26:8).
%	}
}

\Verse{10}
{
	\hE{וַתֹּאמֶר לְאַבְרָהָם גָּרֵשׁ הָאָמָה הַזֹּאת וְאֶת בְּנָהּ כִּי לֹא יִירַשׁ בֶּן הָאָמָה הַזֹּאת עִם בְּנִי עִם יִצְחָק׃}
}
{
	\eE{
		And she said to Abraham, “Drive this maid out—and her son,
		because the son of this maid will not inherit with my son,
		with Isaac.”
	}
}
{
	So Sarah's like ``Remember how I said you should get the maid pregnant so you did and but then some YHWH magic happened and so now I have a kid?''

	Abe's listening.

	She continues ``Anyways you should kick her and your other son out because when you die I'm not splitting all the stuff we got from our monarch entrapment side project with her. I mean it's only reasonable, it was ME that had to f---''

	Abe's listening.

	``...f--- ind a way to avoid having to sleep with all those dirty handsome foreign kings.''
}

\Verse{11}
{
	\hE{וַיֵּרַע הַדָּבָר מְאֹד בְּעֵינֵי אַבְרָהָם עַל אוֹדֹת בְּנוֹ׃}
}
{
	\eE{
		And the thing was very bad in Abraham’s eyes in regard to
		his son.
	}
}
{
	Abe's comments on the matter have been withheld for profanity.

%	\footnote{
%		the thing was very bad in Abraham’s eyes in regard to his son. The
%		Israelite who wrote this text—which favors his own ancestor, Isaac—still
%		expressed sympathy for Ishmael. He wrote of Abraham’s affection for him
%		and God’s promise for him.
%	}
}

\Verse{12}
{
	\hE{וַיֹּאמֶר אֱלֹהִים אֶל אַבְרָהָם אַל יֵרַע בְּעֵינֶיךָ עַל הַנַּעַר וְעַל אֲמָתֶךָ כֹּל אֲשֶׁר תֹּאמַר אֵלֶיךָ שָׂרָה שְׁמַע בְּקֹלָהּ כִּי בְיִצְחָק יִקָּרֵא לְךָ זָרַע׃}
}
{
	\eE{
		And God said to Abraham, “Let it not be bad in your eyes
		about the boy and about your maid. Everything that Sarah
		tells you: listen to her voice; because seed for you will
		be called by Isaac.
	}
}
{
	So God says ``Abe your emotions are all wrong. I know you see this `sending your firstborn son and second wife away' situation as a bad thing. Let's not feel that way, eh? Listen to your wife Abe. Your cu--- You're communicating all wrong. Marriage is about comeyounakeding. Nothing. Hear what? What do you mean Ishmael's crying? Abe don't listen to Ishmael. I'll listen to Ishmael. Focus Abe. Listen to me. You listening? Ok, here's my advice. Listen to Sarah. I mean she's your wife Abe. It's not unreasonable for her to ask you to send your wife and son away. You have a wife and son for God's sake Abe it's time to settle down. Sure we all like to sow our wild c--- oats, um, when we're young. Where was I going with this? Oh right. Abe look it's ok to run around wearing wild coats when we're young, but it's time to grow up. You're an adult now. You're a hundred. And one wife is enough Abe. You're a dad now. It's time to settle down and stop running around trying to spread your seed all over the place. You have responsibilities. Which is more important Abe? Being a good father to your wife and son, or running around trying to spread your seed all over the place? So listen to your wife Abe. Be a good dad. It's time to focus on your wife and son. So send away your wife and son. See? Problem solved. Anyways Abe you don't have to listen to me. But listen to your wife. And so do whatever she says. It's time to stop taking care of yourself and start taking care of her. No not her. Don't let me down Abe. I've got big plans for you. We're gonna spread your seed far and wide, to all the ladies of the earth Abe. But not that one. The other one. Right, the one you have a kid with. No not her. God works in mysterious ways Abe. Life's funny sometimes. So don't forget you're my \#1 guy Abe. And I'm gonna spread your seed to all the nations Abe. But your seed's gonna be funny. That's why we added that `ha' to your name. Get it? So sack up and turn that frown upside down and tell your wife (not her, the other one, the quiet nice one) to pack up her sack, tell that old hag to pack her bag and get out of town. And after that, let's get down to business and get you some {\pussya}.

%	\footnote{
%		seed for you will be called by Isaac. Meaning: Abraham’s best-known line
%		will later be identified as children of Abraham-Isaac-Jacob, rather than
%		Abraham-Ishmael or Abraham-Midian or Abraham-Zimran.
%	}
}

\Verse{13}
{
	\hE{וְגַם אֶת בֶּן הָאָמָה לְגוֹי אֲשִׂימֶנּוּ כִּי זַרְעֲךָ הוּא׃}
}
{
	\eE{
		And I’ll make the maid’s son into a nation as well,
		because he’s your seed.”
	}
}
{
	Don't worry Abe we're gonna give him as much {\pussyb} as I'm giving you. He's safe with me! Wave goodbye Abe! Bon voyage!

%	\footnote{
%		I’ll make the maid’s son into a nation as well. Namely, the Ishmaelites.
%		Their genealogy appears in Gen 25:13–18, and they will figure in the
%		Joseph story (37:25–28; 39:1). They do not appear in the Torah after
%		that.
%	}
}

\Verse{14}
{
	\hE{וַיַּשְׁכֵּם אַבְרָהָם בַּבֹּקֶר וַיִּקַּח לֶחֶם וְחֵמַת מַיִם וַיִּתֵּן אֶל הָגָר שָׂם עַל שִׁכְמָהּ וְאֶת הַיֶּלֶד וַיְשַׁלְּחֶהָ וַתֵּלֶךְ וַתֵּתַע בְּמִדְבַּר בְּאֵר שָׁבַע׃}
}
{
	\eE{
		And Abraham got up early in the morning and took bread and
		a bottle of water and gave to Hagar—he put them on her
		shoulder—and the boy and he sent her off. And she went and
		strayed in the Beer-sheba wilderness.
	}
}
{
	So Abe sends his wife and firstborn away.

	Even gets up early in the morning.

	He's a good guy though, so he doesn't send his wife and son away empty handed.

	He puts some bread and water on her shoulder and sends them off.

	Look nobody's perfect.

	Go try to raise a family.

	Even with YHWH following you around, it's not easy.
}

\Verse{15}
{
	\hE{וַיִּכְלוּ הַמַּיִם מִן הַחֵמֶת וַתַּשְׁלֵךְ אֶת הַיֶּלֶד תַּחַת אַחַד הַשִּׂיחִם׃}
}
{
	\eE{
		And the water was finished from the bottle, and she thrust
		the boy under one of the shrubs,
	}
}
{
	Once Hagar finishes the bottle of water, she throws the baby under a bush.

	Also note the baby is a 14 year old guy at this point.
}

\Verse{16}
{
	\hE{וַתֵּלֶךְ וַתֵּשֶׁב לָהּ מִנֶּגֶד הַרְחֵק כִּמְטַחֲוֵי קֶשֶׁת כִּי אָמְרָה אַל אֶרְאֶה בְּמוֹת הַיָּלֶד וַתֵּשֶׁב מִנֶּגֶד וַתִּשָּׂא אֶת קֹלָהּ וַתֵּבְךְּ׃}
}
{
	\eE{
		and she went and sat opposite, going as far as a bow-shot,
		because she said, “Let me not see the boy’s death.” And
		she sat opposite and raised her voice and wept.
	}
}
{
	She goes and sits across from the bush, about as far as you can shoot an arrow.

	Then she cries about how she doesn't want to see her son die.

	Meanwhile, I imagine Ishmael is confused.

	I picture him standing up and slowly and awkwardly walking away from the bush.

	He doesn't need that bush. He's 14. He's got a bush of his own.
}

\Verse{17}
{
	\hE{וַיִּשְׁמַע אֱלֹהִים אֶת קוֹל הַנַּעַר וַיִּקְרָא מַלְאַךְ אֱלֹהִים אֶל הָגָר מִן הַשָּׁמַיִם וַיֹּאמֶר לָהּ מַה לָּךְ הָגָר אַל תִּירְאִי כִּי שָׁמַע אֱלֹהִים אֶל קוֹל הַנַּעַר בַּאֲשֶׁר הוּא שָׁם׃}
}
{
	\eE{
		And God heard the boy’s voice. And an angel of God called
		to Hagar from the heavens and said to her, “What trouble
		do you have, Hagar? Don’t be afraid. Because God has heard
		the boy’s voice where he is.
	}
}
{
	So God messages Hagar.
	
	He's like ``What's wrong now, Hagar? Don't be afraid. I hear Ishmael. Get it?''

	So both the Narrator and God both make the ``Ishmael'' joke in this line. Again.
}

\Verse{18}
{
	\hE{קוּמִי שְׂאִי אֶת הַנַּעַר וְהַחֲזִיקִי אֶת יָדֵךְ בּוֹ כִּי לְגוֹי גָּדוֹל אֲשִׂימֶנּוּ׃}
}
{
	\eE{
		Get up. Carry the boy and hold him up in your hand,
		because I shall make him into a big nation.”
	}
}
{
	God says to pick the boy up and hold him in her hand.

	Don't worry, it makes sense.

	He's going to be a country.

%	\footnote{
%		I shall make him into a big nation. These words are spoken by the angel,
%		but they certainly appear to be the words of God, not the angel itself.
%		Angels do not make great nations in the Tanak. This is another
%		demonstration that angels are not independent beings in the Hebrew Bible
%		but are rather expressions of God’s presence. (See the comment on Gen
%		Footnote 18:3.)
%	}
}

\Verse{19}
{
	\hE{וַיִּפְקַח אֱלֹהִים אֶת עֵינֶיהָ וַתֵּרֶא בְּאֵר מָיִם וַתֵּלֶךְ וַתְּמַלֵּא אֶת הַחֵמֶת מַיִם וַתַּשְׁקְ אֶת הַנָּעַר׃}
}
{
	\eE{
		And God opened her eyes, and she saw a water well, and she
		went and filled the bottle with water and had the boy
		drink.
	}
}
{
	So God does a miracle where he opens her eyes for her and then she can see that there's water nearby.

	Now her eyes are opened, so she fills up her water bottle and has the boy drink.
	
	He needs help drinking, he's just a high school freshman.
}

\Verse{20}
{
	\hE{וַיְהִי אֱלֹהִים אֶת הַנַּעַר וַיִּגְדָּל וַיֵּשֶׁב בַּמִּדְבָּר וַיְהִי רֹבֶה קַשָּׁת׃}
}
{
	\eE{
		And God was with the boy, and he grew, and he lived in the
		wilderness and was a bowman.
	}
}
{
	Then the boy grows.

	And he lives in the woods and does archery.
}

\Verse{21}
{
	\hE{וַיֵּשֶׁב בְּמִדְבַּר פָּארָן וַתִּקַּח לוֹ אִמּוֹ אִשָּׁה מֵאֶרֶץ מִצְרָיִם׃}
}
{
	\eE{
		And he lived in the Paran wilderness, and his mother took
		him a wife from the land of Egypt.
	}
}
{
	He lives in the Place Of Caves wilderness.

	His mom visits home and brings him a wife from Egypt.
}

\Verse{22}
{
	\hE{וַיְהִי בָּעֵת הַהִוא וַיֹּאמֶר אֲבִימֶלֶךְ וּפִיכֹל שַׂר צְבָאוֹ אֶל אַבְרָהָם לֵאמֹר אֱלֹהִים עִמְּךָ בְּכֹל אֲשֶׁר אַתָּה עֹשֶׂה׃}
}
{
	\eE{
		And it was at that time, and Abimelek and Phichol, the
		commander of his army, said to Abraham, saying, “God is
		with you in everything that you do.
	}
}
{
	It's that time.

	So My Father Is King the king and Mouth Of All the army commander say to Abe ``You and God always do everything together.''
}

\Verse{23}
{
	\hE{וְעַתָּה הִשָּׁבְעָה לִּי בֵאלֹהִים הֵנָּה אִם תִּשְׁקֹר לִי וּלְנִינִי וּלְנֶכְדִּי כַּחֶסֶד אֲשֶׁר עָשִׂיתִי עִמְּךָ תַּעֲשֶׂה עִמָּדִי וְעִם הָאָרֶץ אֲשֶׁר גַּרְתָּה בָּהּ׃}
}
{
	\eE{
		And now, swear to me here by God that you won’t act
		falsely to me and to my offspring and to my posterity Like
		the kindness that I’ve done with you, you’ll do with me
		and with the land in which you’ve resided.”
	}
}
{
	Promise me you won't lie or cheat me and my kids and their kids. Like I've been nice to you, you know. Be nice to me too. Also be nice to the land. Promise me that, and make it a God promise.
}

\Verse{24}
{
	\hE{וַיֹּאמֶר אַבְרָהָם אָנֹכִי אִשָּׁבֵעַ׃}
}
{
	\eE{And Abraham said, “I’ll swear,”}
}
{
	Abe promises.
}

\Verse{25}
{
	\hE{וְהוֹכִחַ אַבְרָהָם אֶת אֲבִימֶלֶךְ עַל אֹדוֹת בְּאֵר הַמַּיִם אֲשֶׁר גָּזְלוּ עַבְדֵי אֲבִימֶלֶךְ׃}
}
{
	\eE{
		and Abraham criticized Abimelek about a water well that
		Abimelek’s servants had seized.
	}
}
{
	Abe immediately starts bitching at Abimelek about a well.
}

\Verse{26}
{
	\hE{וַיֹּאמֶר אֲבִימֶלֶךְ לֹא יָדַעְתִּי מִי עָשָׂה אֶת הַדָּבָר הַזֶּה וְגַם אַתָּה לֹא הִגַּדְתָּ לִּי וְגַם אָנֹכִי לֹא שָׁמַעְתִּי בִּלְתִּי הַיּוֹם׃}
}
{
	\eE{
		And Abimelek said, “I don’t know who did this thing; and
		also, you didn’t tell me; and also, I didn’t hear of it
		except for today.”
	}
}
{
	And Abimelek is like ``Abe, I'm the king. Do you have any idea how many guys I'm in charge of? I have no idea what any of them are doing. This is the first I've heard about this well situation.''
}

\Verse{27}
{
	\hE{וַיִּקַּח אַבְרָהָם צֹאן וּבָקָר וַיִּתֵּן לַאֲבִימֶלֶךְ וַיִּכְרְתוּ שְׁנֵיהֶם בְּרִית׃}
}
{
	\eE{
		And Abraham took sheep and oxen and gave to Abimelek, and
		the two of them made a covenant.
	}
}
{
	So Abe gives him some livestock and the two call it even.

	They make a deal, so they won't have these misunderstandings in the future.
}

\Verse{28}
{
	\hE{וַיַּצֵּב אַבְרָהָם אֶת שֶׁבַע כִּבְשֹׂת הַצֹּאן לְבַדְּהֶן׃}
}
{
	\eE{And Abraham stood seven ewes of the sheep by themselves.}
}
{
	Abe takes seven of the girl sheep and stands them up by themselves.
}

\Verse{29}
{
	\hE{וַיֹּאמֶר אֲבִימֶלֶךְ אֶל אַבְרָהָם מָה הֵנָּה שֶׁבַע כְּבָשֹׂת הָאֵלֶּה אֲשֶׁר הִצַּבְתָּ לְבַדָּנָה׃}
}
{
	\eE{
		And Abimelek said to Abraham, “What are these seven ewes
		that you’ve stood by themselves?”
	}
}
{
	Abimelek's like ``Whatcha doing?''
}

\Verse{30}
{
	\hE{וַיֹּאמֶר כִּי אֶת שֶׁבַע כְּבָשֹׂת תִּקַּח מִיָּדִי בַּעֲבוּר תִּהְיֶה לִּי לְעֵדָה כִּי חָפַרְתִּי אֶת הַבְּאֵר הַזֹּאת׃}
}
{
	\eE{
		And he said, “Because you’ll take these seven ewes from my
		hand so that it will be evidence for me that I dug this
		well.”
	}
}
{
	Abe responds ``You'll take these girl sheep. That'll prove I dug the well.''
}

\Verse{31}
{
	\hE{עַל כֵּן קָרָא לַמָּקוֹם הַהוּא בְּאֵר שָׁבַע כִּי שָׁם נִשְׁבְּעוּ שְׁנֵיהֶם׃}
}
{
	\eE{
		On account of this he called that place Beer-sheba,
		because the two of them swore there
	}
}
{
	On account of this he called the place Beer sheba, for Abe was drunk.

	And doing something with seven%
	\footnote{Lit. še\heb{ב}a\heb{א} in the original text.}
	female sheep.%
	\footnote{
		Lit. she--- she\korean{ㅂ}as in the original text.
	}
}

\Verse{32}
{
	\hE{וַיִּכְרְתוּ בְרִית בִּבְאֵר שָׁבַע וַיָּקם אֲבִימֶלֶךְ וּפִיכֹל שַׂר צְבָאוֹ וַיָּשֻׁבוּ אֶל אֶרֶץ פְּלִשְׁתִּים׃}
}
{
	\eE{
		and made a covenant in Beer-sheba. And Abimelek and
		Phichol, the commander of his army, got up and went to the
		land of the Philistines.
	}
}
{
	So that's the deal.

	Then My Father Is King the king, and Mouth Of All, the army commander, got up and went to Palestine.
}

\Verse{33}
{
	\hE{וַיִּטַּע אֶשֶׁל בִּבְאֵר שָׁבַע וַיִּקְרָא שָׁם בְּשֵׁם}
	\hRed{יהוה}
	\hE{אֵל עוֹלָם׃}
}
{
	\eE{
		And he planted a tamarisk in Beer-sheba and he called the
		name of
	}
	\eRed{YHWH}
	\eE{El Olam.}
}
{
	And he plants a tree in Beer-sheba.

	And he called the name of \eRed{YHWH}%
	\eRed{\footnote{
		This is the third of five unexpected appearances of the name YHWH in a source outside J before Exodus 3.
	}}
	God Forever.
}

\Verse{34}
{
	\hE{וַיָּגר אַבְרָהָם בְּאֶרֶץ פְּלִשְׁתִּים יָמִים רַבִּים׃}
}
{
	\eE{
		And Abraham resided in the land of the Philistines many
		days.
	}
}
{
	And Abe lived in Palestine for a while.
}
